% =====================================================================
%  packages.tex —— 宏包加载与引擎设定
%  引擎：XeLaTeX（中文经 ctex，字体走 Windows 字体集）
%  加载顺序有讲究：biblatex -> hyperref -> cleveref -> imakeidx
%  请勿在此文件写正文，也不要写版式参数（那些属于 style.tex）
% =====================================================================

% ---- 页面几何 --------------------------------------------------------
\usepackage[
  top=25mm, bottom=25mm,
  left=28mm, right=25mm,
  headsep=6mm, footskip=10mm
]{geometry}

% ---- 数学 ------------------------------------------------------------
\usepackage{amsmath}
\usepackage{amssymb}
\usepackage{mathtools}
\usepackage{bm}

% ---- 表格与浮动体 ----------------------------------------------------
\usepackage{booktabs}
\usepackage{array}
\usepackage{tabularx}
\usepackage{longtable}
\usepackage{multirow}
\usepackage{float}
\usepackage{caption}

% ---- 列表 ------------------------------------------------------------
\usepackage{enumitem}

% ---- 可断行的等宽标识符 ----------------------------------------------
%  xurl 让 \path 可以在任意标点处断行。本文档的 \pid（模块标识，
%  形如 physics.quantum.analytic_tw）用它来避免窄表格列里的溢出盒。
%  ⚠ 只给 \pid 用：\path 是 verbatim 扫描，会\emph{吞掉空格}，
%  因此含空格的 \akey/\fpath 继续用 \detokenize（见 style.tex）。
\usepackage{xurl}

% ---- 颜色与图形 ------------------------------------------------------
\usepackage{xcolor}
\usepackage{graphicx}

%  调色板必须在 hyperref 之前定义：hyperref 的 colorlinks 会直接引用这些名字。
\definecolor{NemoLink}{HTML}{1F4E79}   % 深蓝：内部链接
\definecolor{NemoCite}{HTML}{7A3E00}   % 赭石：文献引用
\definecolor{NemoUrl}{HTML}{0B6E4F}    % 深绿：外部链接
\definecolor{NemoRule}{HTML}{4A4A4A}   % 灰：表格线
\definecolor{NemoCardBg}{HTML}{F2F5F8} % 浅蓝灰：模块卡片底纹
\definecolor{NemoKeyBg}{HTML}{EDEDED}  % 浅灰：代码/键名底纹
\definecolor{NemoWarn}{HTML}{9C2A00}   % 暗红：边界与风险提示

% ---- 单位 ------------------------------------------------------------
\usepackage{siunitx}
\sisetup{
  detect-all,
  per-mode = symbol,
  range-phrase = {\text{--}},
  range-units = single,
  separate-uncertainty = true
}

% ---- 参考文献（全局唯一一份，样式为国标 GB/T 7714-2015）--------------
\usepackage[backend=biber,style=gb7714-2015,sorting=gb7714-2015]{biblatex}
\addbibresource{refs.bib}

% ---- 超链接 ----------------------------------------------------------
\usepackage[
  unicode=true,
  bookmarksnumbered=true,
  bookmarksopen=true,
  colorlinks=true,
  linkcolor=NemoLink,
  citecolor=NemoCite,
  urlcolor=NemoUrl,
  filecolor=NemoUrl,
  pdftitle={Huntian: GaN 射频 HEMT 轻量化 TCAD 内核},
  pdfsubject={模块化技术文档},
  pdfkeywords={GaN, HEMT, TCAD, Schroedinger-Poisson, drift-diffusion}
]{hyperref}

% ---- 智能交叉引用 ----------------------------------------------------
%  注意：这里\emph{不}加载 cleveref —— 它与 ctex 不兼容，第二遍编译会报
%  "Extra \endcsname"（最小复现已确认，且与 hyperref 无关）。
%  本文档改用自己的引用宏，见 preamble/refs.tex。

% ---- 索引 ------------------------------------------------------------
%  P0 用 texindy（UTF-8 + 英文排序规则）；中文条目一律带 ASCII 排序键
%  （写成 \index{3method@方法!...}），因此排序完全由文档控制，
%  与 makeindex / xindy 的选择无关。详见 §"索引与术语约定"。
\usepackage{imakeidx}
\makeindex[
  title={索引},
  intoc=true,
  program=texindy,
  options={-C utf8 -L english},
  columns=2,
  columnsep=1.2em
]
